\section{Segmento}

O segmento $[a, b]$ vai de $a$ até $b$, incluindo as pontas.

\subsection*{Ponto no segmento}

\begin{drawing}
  \fill[fig_fill] (2,0) circle (2);
  \draw[fig_accent, dashed] (2,0) circle (2);
  \draw[thick] (0,0) -- (4,0);
  \fill (0,0) circle (1.5pt) node[left] {$a$};
  \fill (4,0) circle (1.5pt) node[right] {$b$};
  \fill (1.4,1.1) circle (1.5pt) node[above] {$p$};
  \draw[gray] (0,0) -- (1.4,1.1) -- (4,0);
  \fill (2.8,0) circle (1.5pt) node[below] {$q$};
\end{drawing}

\texttt{in\_disk(a, b, p)} diz se $p$ está no círculo de diâmetro $ab$: o ângulo $a\,p\,b$ é de pelo menos $90^\circ$,
ou seja, \texttt{dot(a - p, b - p)} $\le 0$. \texttt{on\_segment(a, b, p)} junta isso com \texttt{orient} $= 0$: o
ponto está na reta e entre as pontas, como $q$.

\subsection*{Interseções}

\begin{drawing}
  \draw[thick] (0,0) -- (2.4,1.8);
  \draw[thick] (0,1.8) -- (2.4,0);
  \fill[fig_highlight] (1.2,0.9) circle (2pt);
  \node at (1.2,-0.6) {proper};
  \draw[thick] (3.4,0) -- (5.8,0);
  \draw[thick] (4.6,0) -- (4.6,1.8);
  \fill[fig_highlight] (4.6,0) circle (2pt);
  \node at (4.6,-0.6) {encosta};
  \draw[thick] (6.8,0.75) -- (8.8,0.75);
  \draw[fig_accent, thick] (7.8,1.05) -- (9.6,1.05);
  \draw[fig_highlight, dotted] (7.8,0.75) -- (7.8,1.05);
  \draw[fig_highlight, dotted] (8.8,0.75) -- (8.8,1.05);
  \fill[fig_highlight] (7.8,0.9) circle (2pt);
  \fill[fig_highlight] (8.8,0.9) circle (2pt);
  \node at (8.2,-0.6) {sobrepostos};
\end{drawing}

\begin{itemize}
  \item \texttt{proper(a, b, c, d)}: os segmentos se cruzam num ponto interior aos dois. Acontece quando $c$ e $d$
        estão em lados opostos da reta $ab$ e $a$ e $b$ em lados opostos da reta $cd$ (sinais de \texttt{orient}
        com produto negativo). Exato com \texttt{T} inteiro.
  \item \texttt{intersects(a, b, c, d)}: os segmentos têm algum ponto em comum, em qualquer um dos três casos do
        desenho. É o \texttt{proper} ou alguma ponta de um em cima do outro (\texttt{on\_segment}). Exato com
        \texttt{T} inteiro.
  \item \texttt{proper\_inter(a, b, c, d)} devolve esse ponto de cruzamento.
  \item \texttt{inters(a, b, c, d)} devolve todos os casos: o ponto do cruzamento, a ponta que encosta no outro
        segmento, ou as duas pontas do trecho em comum quando são colineares e sobrepostos.
\end{itemize}

\subsection*{Distâncias}

\begin{drawing}
  \fill[fig_fill] (0,-1.2) rectangle (4,1.8);
  \draw[gray, dashed] (0,-1.2) -- (0,1.8);
  \draw[gray, dashed] (4,-1.2) -- (4,1.8);
  \draw[thick] (0,0) -- (4,0);
  \fill (0,0) circle (1.5pt) node[above left] {$a$};
  \fill (4,0) circle (1.5pt) node[below right] {$b$};
  \fill (2.2,1.4) circle (1.5pt) node[above] {$p$};
  \draw[fig_highlight, thick] (2.2,1.4) -- (2.2,0);
  \fill (5.6,1.2) circle (1.5pt) node[above] {$p'$};
  \draw[fig_highlight, thick] (5.6,1.2) -- (4,0);
  \fill (-1.4,-0.9) circle (1.5pt) node[left] {$p''$};
  \draw[fig_highlight, thick] (-1.4,-0.9) -- (0,0);
\end{drawing}

\texttt{segment\_point(a, b, p)}: se $p$ está na faixa azul (à frente de $a$ e de $b$, olhando para dentro do
segmento), a distância é até a reta, pela perpendicular; fora dela, é até a ponta mais perto. \\
\texttt{segment\_segment(a, b, c, d)} é $0$ se os segmentos se cruzam; senão, os pontos mais próximos incluem uma
das quatro pontas, então é o menor entre as distâncias de cada ponta até o outro segmento.

\subsection*{Distância até uma semirreta}

\begin{drawing}
  \draw[->, thick] (0,0) -- (3.5,0);
  \fill (0,0) circle (1.5pt) node[below] {$c$};
  \draw[fig_accent, thick] (1,-1.2) -- (2.6,1.0);
  \fill[fig_highlight] (1.873,0) circle (2pt);
  \node at (1.6,-1.8) {se cruzam: 0};
  \draw[->, thick] (5.5,0) -- (9,0);
  \fill (5.5,0) circle (1.5pt) node[below] {$c$};
  \draw[fig_accent, thick] (6.8,1.6) -- (8.4,0.9);
  \draw[fig_highlight, thick] (8.4,0.9) -- (8.4,0);
  \draw[fig_accent, thick] (4.2,-1.5) -- (4.6,0.9);
  \draw[fig_highlight, thick] (5.5,0) -- (4.47,0.17);
  \node at (7,-1.8) {ponta do segmento ou origem $c$};
\end{drawing}

\texttt{segment\_ray(a, b, c, d)} é a distância do segmento $[a, b]$ até a semirreta \texttt{(c, d)}. Usa o mesmo
teste do \texttt{ray\_ray}, com o segmento escrito como $a + s \cdot w$ e $0 \le s \le 1$: além dos sinais, testa
\texttt{cross(a - c, v)} $\le$ \texttt{cross(v, w)} (comparando com o sinal do denominador). Se não se cruzam, a
resposta é a menor distância de uma ponta $a$ ou $b$ até a semirreta, ou da origem $c$ até o segmento.

\textbf{Exato ou real.} \texttt{in\_disk}, \texttt{on\_segment}, \texttt{proper}, \texttt{intersects} e as distâncias (inclusive
\texttt{segment\_ray}) funcionam com
\texttt{T} inteiro. \texttt{proper\_inter} e \texttt{inters} devolvem pontos novos e exigem \texttt{T} real.
